% Lecture slides for Guide 2, "Tokenization: how text becomes numbers".
% Written from content/02-tokenizer.mdx. Code, outputs, figures and notebook cell numbers
% are generated (tools/build_slide_assets.py, tools/export-figures.mjs); every other number
% must appear in the guide, which tools/check_slides.py checks. Build: node tools/build-slides.mjs
\providecommand\primeropts{}
\documentclass[aspectratio=169,11pt,t,\primeropts]{beamer}
\def\guidename{02-tokenizer}
\input{generated/02-tokenizer/meta.tex}
\input{generated/02-tokenizer/cells.tex}
\usetheme{primer}

\title{Tokenization}
\subtitle{How text becomes numbers}

% a token as printed: monospace, a visible space where the token holds one
\newcommand\tok[1]{\texttt{#1}}
\newcommand\tb[1]{{\setlength\fboxsep{1.5pt}\fcolorbox{primerRuleStrong}{white}{\texttt{#1}}}}% one boxed token
\newcommand\vsp{\textvisiblespace}% U+2423, from the DejaVu fallback

\begin{document}

% ======================================================================== opening
\begin{frame}[plain,noframenumbering]
  \titlepage
\end{frame}

\begin{frame}{Where we are}
  \begin{center}\primerseriesmap\end{center}
  \vspace{2mm}
  \begin{itemize}
    \item \recall{1}{our toy network read one character at a time: a one-hot row of 27 numbers
      in, a probability for each of 27 next characters out.}
    \item Real language models don't read characters. Today: what they read instead, and why.
  \end{itemize}
\end{frame}

\begin{frame}{Why an engineer should care}
  As of October 2026, the major hosted LLM APIs \textbf{charge by the token} and \textbf{cap
  their input in tokens}.
  \vspace{4mm}
  \begin{columns}[T]
    \column{0.31\textwidth}
    \begin{keyidea}[Length]
      how many tokens a text becomes sets what it costs and whether it fits
    \end{keyidea}
    \column{0.31\textwidth}
    \begin{keyidea}[Vocabulary]
      the list of tokens sets the size of the model's output
    \end{keyidea}
    \column{0.31\textwidth}
    \begin{keyidea}[What the model sees]
      only integers: spelling, digits and other languages arrive pre-chopped
    \end{keyidea}
  \end{columns}
  \vspace{4mm}
  Guide 1 asked you to think of a token as ``roughly a word''. That was a placeholder; today
  repays it.
\end{frame}

\begin{frame}{Learning objectives}
  By the end of this lecture you'll be able to:
  \begin{enumerate}
    \item \textbf{say} what a token, a token ID and a tokenizer are, and why the unit is neither
      a character nor a word;
    \item \textbf{build} a small tokenizer's vocabulary by hand with byte-pair encoding;
    \item \textbf{encode} text the tokenizer has never seen, and \textbf{explain} why bytes mean
      nothing is ever unknown;
    \item \textbf{run} GPT-2's real tokenizer and \textbf{read} its output;
    \item \textbf{predict} where tokenization makes a model stumble: spelling, digits and text
      that isn't English.
  \end{enumerate}
\end{frame}

% ======================================================================== section 1
\section{The unit a model reads}

\begin{frame}{Text has to become numbers}
  \begin{itemize}
    \item A network computes with numbers, so text has to become numbers first.
    \item Guide 1 did it the simplest way: a fixed list of 27 characters (the space and a to z),
      each character standing for its index in the list.
    \item A language model does the same thing with a \textbf{longer list} and \textbf{bigger
      units}.
  \end{itemize}
  \pause
  \vspace{2mm}
  \begin{itemize}
    \item One unit of text is a \term{token}. The list is the \textbf{vocabulary} $\mathcal{V}$,
      of size $\lvert\mathcal{V}\rvert$ (27 in Guide 1). The model reads from the same list it
      writes to.
    \item A token's index in the list is its \term{token ID}, an integer counted from 0.
  \end{itemize}
\end{frame}

\begin{frame}{The tokenizer converts in both directions}
  \begin{center}
  \begin{tikzpicture}[node distance=1.2cm,
      box/.style={draw=primerRuleStrong,rounded corners=2pt,minimum height=0.75cm,font=\small,inner sep=4pt},
      net/.style={box,fill=primerAccent!10,draw=primerAccent}]
    \node[box] (txt) {text};
    \node[box,right=2.2cm of txt] (ids) {token IDs};
    \node[net,right=1.8cm of ids] (nn) {network};
    \draw[-{Stealth},primerInk] ([yshift=1.5mm]txt.east) -- node[above,font=\small] {encode} ([yshift=1.5mm]ids.west);
    \draw[-{Stealth},primerInk] ([yshift=-1.5mm]ids.west) -- node[below,font=\small] {decode} ([yshift=-1.5mm]txt.east);
    \draw[-{Stealth},primerInk] (ids) -- (nn);
    \draw[dashed,primerMuted] ($(ids.north west)+(-0.5,0.5)$) rectangle ($(nn.south east)+(0.3,-0.6)$);
    \node[font=\scriptsize,primerMuted,anchor=north east] at ($(nn.south east)+(0.3,-0.6)$) {the ID side: no letter ever gets here};
  \end{tikzpicture}
  \end{center}
  \vspace{2mm}
  \begin{itemize}
    \item The \term{tokenizer} holds the list. It \term{encode}s text into a sequence of token IDs
      and \term{decode}s token IDs back into text.
    \item The network receives integers and outputs a probability for each integer in the list.
  \end{itemize}
  \pause
  \begin{keyidea}[The question for today]
    What should the unit be: a character, a word, or something in between? The choice sets how
    long every input is and how big the list has to be.
  \end{keyidea}
\end{frame}

\begin{frame}{Check your understanding}
  \checkq{Where does the tokenizer sit relative to the network, and what crosses the boundary
    between them?}
  {The tokenizer encodes text into token IDs before the network, and decodes token IDs back into
    text. The network itself only receives integers and outputs a probability for each one; it
    never sees a letter.}
\end{frame}

% ======================================================================== section 2
\section{Two bad extremes}

\begin{frame}{Count Guide 1's toy text two ways}
  \primercode{27}
  \primeroutput{27}
  \cellnote{27}
  \vspace{1mm}
  {\small The code is plain Python until GPT-2, in blocks that run in order and share variables,
  like notebook cells.\par}
\end{frame}

\begin{frame}{Characters: a tiny list, long sequences}
  \begin{center}\small
  \begin{tabular}{@{}lcc@{}}
    \toprule
    \textbf{Unit} & \textbf{Tokens for the toy text} & \textbf{Distinct tokens (vocabulary)}\\
    \midrule
    Characters & 128 & 15\\
    Words & 33 & 12\\
    \bottomrule
  \end{tabular}
  \end{center}
  \begin{itemize}
    \item Only 15 of Guide 1's 27 characters turn up; the text never uses the other 12 letters.
    \item<2-> \textbf{Characters}: the vocabulary is tiny, and a new word can't catch it out, since
      every new word is spelled with old letters.
    \item<3-> \textbf{The price is length}: 128 tokens to say 33 words, and one token carries very
      little. A model spends computation on every token it reads, and Guide 4 shows how fast that
      cost grows with length.
  \end{itemize}
\end{frame}

\begin{frame}{Words: short sequences, until a new word arrives}
  \primercode{51}
  \primeroutput{51}
  \cellnote{51}
\end{frame}

\begin{frame}{The rug is gone}
  \begin{itemize}
    \item ``rug'' has no ID: the tokenizer falls back on the \term{unknown token}, \tok{<unk>}, ID 12.
    \item Encoding no longer fails, but decoding returned \tok{the cat sat on the <unk>}: every
      unseen word becomes the same ID 12, and nothing after the tokenizer can tell which it was.
  \end{itemize}
  \pause
  \begin{itemize}
    \item A longer list has the same hole: real text keeps producing names, typos, product codes,
      slang and variable names. A word vocabulary grows without limit or stays full of gaps.
    \item It also wastes what it has: ``cat'' and ``cats'' would be two unrelated entries.
  \end{itemize}
  \pause
  \begin{keyidea}[What we want, all at once]
    Short sequences, no unknown tokens, a vocabulary of manageable size. Two moves: pieces of
    words chosen by counting, with bytes underneath as a safety net.
  \end{keyidea}
\end{frame}

\begin{frame}{One sentence, three splits}
  \begin{columns}[c]
    \column{0.66\textwidth}
    \vspace{-3mm}
    \primerfigure[0.8\textheight]{1}
    \column{0.32\textwidth}
    \begin{itemize}
      \item \textbf{Characters}: 22 tokens.
      \item \textbf{Words}: 6 tokens.
      \item \textbf{The tokenizer we build next}: 11 tokens, pieces of words.
    \end{itemize}
    \vspace{2mm}
    {\footnotesize\color{primerMuted}A small bracket marks a space. Figure 1 of the guide.\par}
  \end{columns}
\end{frame}

\begin{frame}{Check your understanding}
  \checkq{The toy text is 128 characters and 33 words. Give one reason to prefer each unit, and
    the failure that rules each one out.}
  {Characters: a tiny vocabulary (15 here), and no new word is unknown as long as it's spelled
    with letters already in the list, but sequences are about
    four times longer, and every token costs computation. Words: short sequences, but any unseen
    word collapses to \tok{<unk>} and is lost, and related forms like ``cat'' and ``cats'' share
    nothing.}
\end{frame}

% ======================================================================== section 3
\section{Learning the vocabulary from data: byte-pair encoding}

\begin{frame}{The idea: glue the most frequent pair, again and again}
  \begin{enumerate}
    \item Start with characters as the tokens.
    \item<2-> Find the pair of adjacent tokens that occurs most often in the text.
    \item<3-> Glue it into one new token and add that token to the vocabulary.
    \item<4-> Repeat.
  \end{enumerate}
  \uncover<5->{%
  \begin{itemize}
    \item Each glued pair is a \term{merge}; the procedure is \term{byte-pair encoding (BPE)}.
    \item The name comes from a compression algorithm that merged pairs of bytes (Gage, 1994),
      adapted for translation models by Sennrich, Haddow \& Birch (2016). We merge characters
      for now; bytes come back two sections on.
  \end{itemize}}
\end{frame}

\begin{frame}{One rule first: pre-tokenization}
  \term{Pre-tokenization}: split the text into words first, and let merges happen only inside a
  word, never across two.
  \begin{itemize}
    \item For now a word is a run of letters (GPT-2's rule is fussier; we'll see it run).
    \item<2-> Without it, a token could straddle two words, useful only when those two words meet
      again. The trainer would also rescan the whole text after every merge.
    \item<3-> With it, no pair spans a word boundary, so the trainer keeps a table of
      \textbf{distinct words and their counts} in place of the full text.
  \end{itemize}
  \uncover<4->{%
  \begin{keyidea}[Where do the spaces go?]
    Decoding must return the original text exactly, so each word keeps the space in front of it:
    \tok{the}, \tok{\vsp cat}, \tok{\vsp sat}, \ldots\ You'll see it in GPT-2's output, where
    most tokens carry a leading space.
  \end{keyidea}}
\end{frame}

\begin{frame}{The trainer, part 1: gluing one pair inside a word}
  \primercodepart{89}{1}
  \vspace{1mm}
  {\small \texttt{merge\_pair} walks along a word (a tuple of tokens) and glues every occurrence of
  the pair \texttt{(a, b)} into one token.\par}
  \cellnote{89}
\end{frame}

\begin{frame}{The trainer, part 2: a table of words, not the whole text}
  \primercodepart{89}{2}
  \vspace{2mm}
  \begin{itemize}
    \item The pattern \texttt{" ?[a-z]+"} is pre-tokenization: a run of letters, with the space in
      front of it if there is one.
    \item \texttt{Counter} maps each distinct word, as a tuple of tokens, to how often it occurs.
      The trainer works on this table from now on.
    \item \texttt{merges} will hold the learned pairs, in order.
  \end{itemize}
\end{frame}

\begin{frame}{The trainer, part 3: count, pick, merge, repeat}
  \primercodepart{89}{3}
\end{frame}

\begin{frame}{The trainer, part 4: ten merges on the toy text}
  \renewcommand\codesize{\footnotesize}
  \primercodepart{89}{4}
  \vspace{2mm}
  {\small Each step of the loop counts every adjacent pair, weighted by how often its word occurs,
  takes the most frequent, merges it everywhere and prints the text's new length.\par}
  \vspace{3mm}
  \begin{keyidea}[Predict first]
    The text starts at 128 tokens. Which pair do you expect to be merged first, and roughly how
    long is the text after 10 merges?
  \end{keyidea}
\end{frame}

\begin{frame}{The trainer's output: the merge table}
  \renewcommand\codesize{\small}
  \primeroutput{89}
  \cellnote{89}
\end{frame}

\begin{frame}{Reading the merge table}
  \begin{itemize}
    \item \tok{t} and \tok{h} sit next to each other 11 times (10 in ``the'', once in ``then''):
      they become \tok{th}, then \tok{th} + \tok{e} becomes \tok{the}.
    \item<2-> Merge 3 attaches the space: seen 10 times, not 11, because the text's very first word
      has no space before it. So \tok{the} and \tok{\vsp the} are two different tokens.
    \item<3-> By merge 8 the algorithm has built \tok{\vsp cat}. Nobody told it which words matter:
      frequent words became single tokens because their letters keep turning up side by side.
    \item<4-> \textbf{Ties.} Merge 1 was a three-way tie (\tok{h}+\tok{e} and \tok{\vsp}+\tok{t} also
      occur 11 times). \texttt{max} takes the pair that sorts last alphabetically: arbitrary, but
      fixed, so your run prints the same table.
  \end{itemize}
\end{frame}

\begin{frame}{The first four merges, step by step}
  \begin{columns}[c]
    \column{0.66\textwidth}
    \vspace{-3mm}
    \primerfigure[0.7\textheight]{2}
    \column{0.32\textwidth}
    \begin{itemize}
      \item Each merge glues one pair of neighbouring tokens wherever it occurs.
      \item The count underneath is the length of the whole toy text after that merge.
      \item ``dog'' is untouched until merges 6 and 10.
    \end{itemize}
    \vspace{2mm}
    {\footnotesize\color{primerMuted}Figure 2 of the guide, on three words of the toy text.\par}
  \end{columns}
\end{frame}

\mode<handout>{%
\begin{frame}{The first four merges, every step (1 of 2)}
  \begin{columns}[T]
    \column{0.49\textwidth}\primerfigonly[0.66\textheight]{2}{1}
    \column{0.49\textwidth}\primerfigonly[0.66\textheight]{2}{2}
  \end{columns}
  {\footnotesize\color{primerMuted}The starting split, then merge 1 (Figure 2's steps 1 and 2).\par}
\end{frame}
\begin{frame}{The first four merges, every step (2 of 2)}
  \begin{columns}[T]
    \column{0.49\textwidth}\primerfigonly[0.66\textheight]{2}{3}
    \column{0.49\textwidth}\primerfigonly[0.66\textheight]{2}{4}
  \end{columns}
  {\footnotesize\color{primerMuted}Merges 2 and 3 (steps 3 and 4); merge 4 is on the previous page.\par}
\end{frame}}

\begin{frame}{The number of merges is the dial}
  After 10 merges: 25 vocabulary entries (the 15 characters plus one per merge), and the toy text
  is 65 tokens, about half of its 128 characters.
  \vspace{2mm}
  \begin{center}
  \begin{tikzpicture}[x=1cm,y=1cm]
    \draw[primerRuleStrong,line width=1pt,{Stealth}-{Stealth}] (0,0) -- (10,0);
    \node[anchor=north,font=\small,align=center] at (0,-0.15) {stop early:\\near characters};
    \node[anchor=north,font=\small,align=center] at (10,-0.15) {no pairs left:\\every word one token};
    \node[anchor=south,font=\small,primerMuted] at (5,0.12) {each extra merge: one more entry, a slightly shorter text};
  \end{tikzpicture}
  \end{center}
  \begin{center}\small GPT-2's designers stopped at 50,000 merges.\end{center}
  \pause
  \begin{itemize}
    \item \textbf{Why stop at all?} \recall{1}{a network outputs one logit per vocabulary entry.}
      Every entry adds parameters and work at the output, and a rare entry turns up too seldom to
      be learned well.
    \item<3-> 50,000 is a chosen trade-off between sequence length and vocabulary size.
  \end{itemize}
\end{frame}

\begin{frame}{Trained before the model exists}
  \begin{itemize}
    \item A tokenizer is trained like this \textbf{once}, on a large sample of text, before the
      model exists.
    \item ``Trained'' here means counting and merging, with \textbf{no gradient descent}.
    \item Its merge list is then frozen, and the model is trained on the token IDs it outputs.
  \end{itemize}
  \vspace{3mm}
  \begin{keyidea}[Two different kinds of learning]
    The tokenizer's vocabulary is fixed by counting, then frozen. The model's parameters are
    learned afterwards, by gradient descent (Guide 1), on the token IDs the tokenizer outputs.
  \end{keyidea}
\end{frame}

\begin{frame}{Check your understanding}
  \checkq{After merge 3 the toy vocabulary contains both \tok{the} and \tok{\vsp the}. Why two
    tokens, and why was the merge seen 10 times rather than 11?}
  {Each word keeps the space in front of it, so ``the'' at the start of the text (no space before
    it) and ``the'' everywhere else are different strings. The pair \tok{\vsp} + \tok{the} occurs
    only where a space precedes ``the'': 10 of its 11 appearances.}
\end{frame}

% ======================================================================== section 4
\section{Encoding text the tokenizer has never seen}

\begin{frame}{Replay the merges, in order}
  Training left an ordered list of merges. To encode new text:
  \begin{enumerate}
    \item pre-tokenize it the same way (accepting any non-space characters, so we can feed it
      letters the trainer never saw);
    \item split each word into characters;
    \item replay the merges \textbf{in the order they were learned}.
  \end{enumerate}
  \pause
  \begin{keyidea}[Why order matters]
    Later merges are built from earlier ones. Merge 7 glues \tok{c} to \tok{at}, and there is no
    \tok{at} until merge 4 has run.
  \end{keyidea}
\end{frame}

\begin{frame}{The encoder}
  \primercodesolo{165}
  \cellnote{165}
\end{frame}

\begin{frame}{The encoder's output}
  \begin{keyidea}[Predict first]
    ``cats'' and ``bat'' never appear in the toy text. Which pieces will each become, and will any
    piece lack an ID?
  \end{keyidea}
  \vspace{2mm}
  \renewcommand\codesize{\small}
  \only<2-| handout:1>{\primeroutput{165}
  \vspace{2mm}
  {\small \texttt{encode} stops at the tokens so you can see them; the IDs come from looking each one
  up in \texttt{vocab}. The last line, 11 tokens, is the third split in Figure 1.\par}}
\end{frame}

\begin{frame}{``cats'' works; ``bat'' doesn't}
  \begin{columns}[T]
    \column{0.48\textwidth}
    \begin{keyidea}[\tok{the cats}]
      ``cats'' appears nowhere in the toy text, and it still encodes cleanly: the known token
      \tok{\vsp cat} followed by the letter \tok{s}. A new word is spelled out of familiar
      pieces, where the word tokenizer flattened it to \tok{<unk>}.
    \end{keyidea}
    \column{0.48\textwidth}
    \uncover<2->{%
    \begin{keyidea}[\tok{the bat}]
      \tok{\vsp}, \tok{b}, \tok{at}: and \tok{b} isn't in the vocabulary, because the toy text has
      no letter b. It has no ID, so this tokenizer would need an unknown token after all.
    \end{keyidea}}
  \end{columns}
  \vspace{2mm}
  \begin{itemize}
    \item Don't read a split as grammar: BPE splits by \textbf{frequency}. A plural common enough
      in the training text gets its own token; only rarer words fall back to shared pieces.
    \item<3-> The vocabulary began as the 15 characters that happened to appear. Capitals, digits,
      punctuation, every other alphabet: still unknown.
  \end{itemize}
\end{frame}

\begin{frame}{Check your understanding}
  \checkq{Why must \texttt{encode} replay the merges in the order they were learned? What would go
    wrong encoding \tok{\vsp cat} if merge 7 (\tok{c} + \tok{at}) ran before merge 4 (\tok{a} +
    \tok{t})?}
  {Merge 7 looks for the token \tok{at}, which only exists after merge 4. Run first, it finds no
    \tok{at} and does nothing, and merge 8 then finds no \tok{cat} to join to the space:
    \tok{\vsp cat} would end as \tok{\vsp}, \tok{c}, \tok{at} instead of the single token
    \tok{\vsp cat}.}
\end{frame}

% ======================================================================== section 5
\section{Bytes: the fallback that never fails}

\begin{frame}{Start from a set that is complete by construction}
  \begin{itemize}
    \item Any text in a file or on a network is stored as a sequence of \term{byte}s, each a number
      from 0 to 255.
    \item \term{UTF-8} (Unicode Transformation Format, 8-bit) is the standard rule for turning
      characters into bytes: one byte for each plain English letter, digit and punctuation mark,
      and between two and four for every other character.
  \end{itemize}
  \vspace{1mm}
  \primercode{196}
  \primeroutput{196}
  \cellnote{196}
  \vspace{1mm}
  {\small ``café'' is 4 characters but 5 bytes (the é takes two: 195 and 169); the cat emoji takes
  four. Whatever the script, the bytes come from the same 256 values.\par}
\end{frame}

\begin{frame}{Byte-level BPE: nothing can be unknown}
  \begin{itemize}
    \item \term{Byte-level BPE} runs the same merging algorithm with the 256 byte values as the
      starting vocabulary.
    \item Frequent byte sequences still merge into pieces of words and whole words.
    \item Text the trainer never saw falls back to shorter pieces and, in the worst case, to single
      bytes, which are \textbf{always} in the vocabulary.
  \end{itemize}
  \pause
  \begin{keyidea}
    A byte-level tokenizer needs no unknown token, because nothing can be unknown to it. Our toy
    started from the characters in the text, which kept its table readable; GPT-2's trainer, like
    many production trainers since, starts from bytes.
  \end{keyidea}
\end{frame}

\begin{frame}{GPT-2's vocabulary, by arithmetic}
  GPT-2 uses byte-level BPE with a vocabulary of 50,257 tokens (Radford et al., 2019):
  \[
    \underbrace{256}_{\text{byte tokens}} \;+\; \underbrace{50{,}000}_{\text{merges}}
    \;+\; \underbrace{1}_{\text{special token}} \;=\; 50{,}257
  \]
  \pause
  \begin{itemize}
    \item A \textbf{special token} is a reserved ID that stands for a control signal and never for a
      piece of text.
    \item GPT-2's is the \term{end-of-text token}, \tok{<|endoftext|>}, ID 50256, the last in the
      list. In a long stream of training text it marks where one document ends and the next
      begins.
    \item GPT-2 has no other special tokens.
  \end{itemize}
\end{frame}

\begin{frame}{Going deeper (optional): other tokenizers you'll meet}
  \small
  \begin{itemize}
    \item \textbf{More special tokens.} Other families mark the beginning or end of a sequence, or
      padding that fills short sequences out to a common length. Tokenizers that don't start from
      bytes also keep an unknown token.
    \item \textbf{WordPiece} (Schuster \& Nakajima, 2012; Wu et al., 2016), BERT's tokenizer
      (Devlin et al., 2019): also merges, but scores a candidate merge by how much more probable it
      makes the training text, and encodes by taking the longest matching piece first. It marks a
      piece that \emph{continues} a word with \tok{\#\#}, where GPT-2 marks a piece that
      \emph{starts} one with a space.
    \item \textbf{Unigram} (Kudo, 2018) works the other way: start from a large vocabulary of
      candidate pieces and prune it down. T5 (Raffel et al., 2020) uses it through the
      SentencePiece library (Kudo \& Richardson, 2018).
    \item \textbf{Normalization.} Some tokenizers lowercase text or strip accents before
      splitting. GPT-2's does neither, so capitals and accents survive encoding.
  \end{itemize}
\end{frame}

\begin{frame}{Check your understanding}
  \checkq{Why does GPT-2's tokenizer need no unknown token, when our toy tokenizer couldn't encode
    ``bat''? And where does 50,257 come from?}
  {It starts from all 256 byte values, and any text is a sequence of bytes, so the worst case is
    single-byte tokens, which are always in the list. The toy started from the 15 characters that
    appeared, so a b had no ID. $256 + 50{,}000 + 1 = 50{,}257$: byte tokens, merges, and the one
    special token.}
\end{frame}

% ======================================================================== section 6
\section{A real tokenizer on a real sentence}

\begin{frame}{The running example meets GPT-2}
  \begin{itemize}
    \item The series traces one short input, ``The cat sat on the'', through every stage of a
      language model. Tokenization is the first stage.
    \item \textbf{Why GPT-2?} Its tokenizer is frozen and small enough to inspect. The model you
      call today has its own tokenizer and will give different counts, so count with that one.
    \item The \texttt{tiktoken} library ships it; outputs come from version 0.14.0, pinned as of
      October 2026. Its first call downloads two vocabulary files (about 1.5 MB in total),
      checks each against a stored hash, and caches them.
  \end{itemize}
\end{frame}

\begin{frame}{Five tokens}
  \primercode{235}
  \primeroutput{235}
  \cellnote{235}
\end{frame}

\begin{frame}{Reading the output}
  \begin{center}\small
  \begin{tabular}{@{}lccccc@{}}
    \toprule
    \textbf{Token} & \tok{The} & \tok{\vsp cat} & \tok{\vsp sat} & \tok{\vsp on} & \tok{\vsp the}\\
    \textbf{ID} & 464 & 3797 & 3332 & 319 & 262\\
    \bottomrule
  \end{tabular}
  \end{center}
  \begin{itemize}
    \item Five tokens, so the \term{sequence length} $n$ is 5.
    \item<2-> Decoding gives back the exact string, capital T included. GPT-2's tokenizer starts
      from bytes, keeps every space inside a token and changes nothing before it splits, so it is
      \term{lossless}: \texttt{decode} always undoes \texttt{encode}.
    \item<3-> The last line confirms the vocabulary size, 50257, and the end-of-text token's ID,
      50256.
  \end{itemize}
\end{frame}

\begin{frame}{The same text, encoded as GPT-2 sees it}
  \vspace{-2mm}
  \primerfigure[0.7\textheight]{3}
  \vspace{-2mm}
  \begin{center}\footnotesize
    Five tokens, four with a leading space (the small bracket), each looked up to get its ID. The
    faded row is the next stage: Guide 3. {\color{primerMuted}(Figure 3 of the guide)}
  \end{center}
\end{frame}

\begin{frame}{Capitals, spaces and the byte fallback}
  \primercodesolo{261}
  \cellnote{261}
\end{frame}

\begin{frame}{Capitals, spaces and the byte fallback: output}
  \begin{keyidea}[Predict first]
    Change the capital or drop the space and you get a different token. Will \tok{\vsp cat},
    \tok{\vsp Cat}, \tok{cat} and \tok{\vsp CAT} share any ID? Will ``Akwirw ier'' need an
    unknown token?
  \end{keyidea}
  \renewcommand\codesize{\footnotesize}
  \only<2-| handout:1>{\primeroutput{261}}
\end{frame}

\begin{frame}{Reading it, line by line}
  \small
  \begin{itemize}
    \item \tok{\vsp cat}, \tok{\vsp Cat}, \tok{cat}, \tok{\vsp CAT}: 3797, 5181, 9246, 38348. To the
      model, \textbf{unrelated integers}. Nothing about 3797 and 5181 says they name the same
      animal; Guide 3 shows how training makes them related.
    \item<2-> ``the bat'', which left the toy tokenizer holding a \tok{b}, is two ordinary tokens.
    \item<3-> The invented ``Akwirw ier'' comes out as six familiar pieces, no unknown token. The
      space stands alone because the merges that build \tok{ier} come earlier in the list and run
      first, and no learned merge joins a space to \tok{ier}, so the space falls back to its own
      single-byte token.
    \item<4-> The emoji's four bytes become three tokens: one merge glued the first two bytes. The
      256 byte tokens fill IDs 0 to 255, in a different order: 238 and 230 are the single-byte
      tokens for bytes 144 and 136.
  \end{itemize}
\end{frame}

\begin{frame}{A simplification, corrected}
  \begin{columns}[T]
    \column{0.48\textwidth}
    \begin{keyidea}[What we taught]
      Pre-tokenization splits the text into words, each keeping its leading space.
    \end{keyidea}
    \column{0.48\textwidth}
    \begin{keyidea}[What GPT-2 does]
      It also separates digits, punctuation and contraction endings such as \tok{'s} into tokens
      of their own, and has its own rule for runs of several spaces.
    \end{keyidea}
  \end{columns}
  \vspace{3mm}
  The apostrophe line shows it: \tok{The}, \tok{\vsp cat}, \tok{'s}, \tok{\vsp 2}, \tok{\vsp mats}.
  \vspace{2mm}
  \begin{itemize}
    \item The same line: \tok{\vsp mats} is one token, as unrelated to \tok{\vsp mat} as
      \tok{\vsp Cat} is to \tok{\vsp cat}, because a frequent plural earns its own entry.
  \end{itemize}
\end{frame}

\begin{frame}{Check your understanding}
  \checkq{\tok{\vsp cat} is ID 3797 and \tok{\vsp Cat} is ID 5181. What does the network know about
    their relationship before training, and why does ``Akwirw ier'' still encode without an
    unknown token?}
  {Nothing: they are two unrelated integers, and only training (Guide 3) relates them. ``Akwirw
    ier'' falls back to shorter pieces the merges do know, and in the worst case to single bytes,
    which are always in the vocabulary.}
\end{frame}

% ======================================================================== section 7
\section{Where tokenization shows through}

\begin{frame}{Predict first: where will the tokenizer show through?}
  \begin{itemize}
    \item How many tokens is ``strawberry'' with no space in front? With one?
    \item Will a seven-digit number split along place value (thousands, hundreds)?
    \item The same sentence in English, Norwegian, German and French: which is cheapest in tokens?
  \end{itemize}
  \vspace{2mm}
  {\small All computed with GPT-2's tokenizer on the next slide.\par}
\end{frame}

\begin{frame}{Spelling, digits and other languages}
  \primercode{299}
  \primeroutput{299}
  \cellnote{299}
\end{frame}

\begin{frame}{Spelling: the model receives IDs, not letters}
  \begin{itemize}
    \item ``strawberry'' with no space before it is three tokens; with one it is a single token.
      Either way, the ten letters aren't among what the model receives.
    \item<2-> It can only learn spelling \textbf{indirectly}: the same word turns up split in
      different ways in the same kinds of sentences, so training can teach it to respond to
      \tok{\vsp strawberry} and to \tok{st}\,+\,\tok{raw}\,+\,\tok{berry} alike. Rarer words split
      all the way down to single letters, linking pieces to letters (Kaushal \& Mahowald, 2022).
    \item<3-> That's why ``how many r's are in strawberry'' became a well-known stumbling block for
      chat models around 2024. Newer models often get it right, but what they receive is still
      IDs.
  \end{itemize}
\end{frame}

\begin{frame}{Numbers: grouped by frequency, not place value}
  \begin{columns}[c]
    \column{0.42\textwidth}
    \begin{center}\large
      \tok{\vsp 1234567} $\to$ \tb{\vsp 123}\,\tb{45}\,\tb{67}\\[3mm]
      \tok{\vsp 2026} $\to$ \tb{\vsp 20}\,\tb{26}\\[2mm]
      {\footnotesize\color{primerMuted}one box per token\par}
    \end{center}
    \column{0.55\textwidth}
    \begin{itemize}
      \item Digits are grouped by which sequences were frequent in the training text, with no
        regard for place value.
      \item<2-> Adding or comparing numbers means lining digits up by place. Here the tokens hold
        three, two and two digits, and another seven-digit number can split differently.
      \item<3-> The model has to work out an alignment the tokens don't hand it: one reason
        arithmetic on long numbers can go wrong.
    \end{itemize}
  \end{columns}
\end{frame}

\begin{frame}{Other languages cost more}
  \begin{center}\small
  \begin{tabular}{@{}lcc@{}}
    \toprule
    \textbf{Sentence} & \textbf{Bytes} & \textbf{GPT-2 tokens}\\
    \midrule
    The cat sat on the mat. & 23 & 7\\
    Katten satt på matten. (Norwegian) & 23 & 9\\
    Die Katze saß auf der Matte. (German) & 29 & 10\\
    Le chat était assis sur le tapis. (French) & 34 & 12\\
    \bottomrule
  \end{tabular}
  \end{center}
  \begin{itemize}
    \item Same meaning; English is cheapest. GPT-2's merges were almost certainly learned from web
      text that was mostly English, so English words tend to be whole tokens.
    \item<2-> The Norwegian sentence has exactly as many bytes as the English one and still needs
      two more tokens.
  \end{itemize}
\end{frame}

\begin{frame}{What you pay for, and what you run out of}
  \begin{itemize}
    \item A model's \term{context length} is the largest number of tokens it can read at once:
      1,024 for GPT-2 (Radford et al., 2019). Models since read far more, but the limit is still a
      token count.
    \item<2-> So the same document \textbf{costs more and fits less} in a language the tokenizer
      splits finely.
    \item<3-> The rule of thumb, about three quarters of a word per token, is folklore tuned to
      English prose. The table's English sentence is 6 words in 7 tokens; the French one, 7 words
      in 12.
  \end{itemize}
  \uncover<4->{%
  \begin{keyidea}[In practice]
    Count with the tokenizer of the model you'll call, on text in the languages you'll send.
  \end{keyidea}}
\end{frame}

\begin{frame}{Check your understanding}
  \checkq{A colleague estimates a French document's API cost from its word count and the
    three-quarters-of-a-word rule. What's wrong with the estimate?}
  {The rule is tuned to English prose. GPT-2 splits other languages into more, smaller pieces: the
    French sentence is 7 words in 12 tokens against English's 6 in 7. The French document is
    likely to cost more and fit less than the estimate says. Count with the model's own tokenizer.}
\end{frame}

% ======================================================================== section 8
\section{From tokens to training examples}

\begin{frame}{The same shift by one position}
  \recall{1}{a training example paired one character with the character that came next.} With
  tokens, the pairing is the same shift by one:
  \vspace{1mm}
  \primercode{337}
  \primeroutput{337}
  \cellnote{337}
  \pause
  \begin{itemize}
    \item Each position's target is the token that follows it: after 464 comes 3797, and after 319
      comes 262.
  \end{itemize}
\end{frame}

\begin{frame}{A probability for every token, never one word}
  \begin{itemize}
    \item What the model sees when it predicts a target is a separate matter: Guide 1's network saw
      one item; Guide 4 gives it the whole preceding sequence.
    \item As in Guide 1, the model outputs a probability for every entry in its list at each
      position: here all 50,257 tokens in $\mathcal{V}$, never one chosen word.
    \item<2-> An ID is still a bare label. Token 3797 isn't larger than token 262 in any sense
      that means something, and a network can't do useful arithmetic on labels.
    \item<3-> For now, an ID is just a label; Guide 3 turns each one into a vector.
  \end{itemize}
  \uncover<4->{%
  \begin{keyidea}
    A tokenizer is a learned vocabulary of frequent pieces of text that turns any text into a
    sequence of integers. It decides how long every input is, what it costs, and which details the
    model gets directly and which it has to learn the hard way.
  \end{keyidea}}
\end{frame}

\begin{frame}{Check your understanding}
  \checkq{For the IDs \texttt{[464, 3797, 3332, 319, 262]}, what are the inputs and targets, and
    why can't the network just use 3797 as a number?}
  {Inputs \texttt{[464, 3797, 3332, 319]}, targets \texttt{[3797, 3332, 319, 262]}. The ID is a
    position in a list, a label: 3797 isn't ``more'' than 262 in any meaningful way, so arithmetic
    on it means nothing until Guide 3 gives each ID a learned vector.}
\end{frame}

% ======================================================================== closing
\section*{Wrap-up}

\begin{frame}{What we simplified, and where it's repaid}
  \begin{center}\small
  \begin{tabular}{@{}L{0.36\textwidth}L{0.42\textwidth}l@{}}
    \toprule
    \textbf{In this guide} & \textbf{In full} & \textbf{Repaid}\\
    \midrule
    \simplification{pre-tokenization splits text into words, each keeping its leading space}{GPT-2 also separates digits, punctuation and contraction endings, and handles runs of spaces}{this guide}
    \addlinespace
    \simplification{the toy BPE starts from the characters in the text}{byte-level BPE starts from the 256 byte values}{this guide}
    \addlinespace
    \simplification{treat a token ID as just a label}{each ID selects a row of a learned table}{Guide 3}
    \addlinespace
    \simplification{the model sees one token at a time}{the whole preceding sequence}{Guide 4}
    \bottomrule
  \end{tabular}
  \end{center}
  \vspace{2mm}
  {\small\color{primerMuted}Repaid today: Guide 1's ``a token is roughly a word''.}
\end{frame}

\begin{frame}{Summary}
  \begin{enumerate}\small
    \item A model reads token IDs. The tokenizer encodes text into IDs and decodes them back; the
      network never sees a letter.
    \item Characters give long sequences; words give unknown tokens. The toy text: 128 characters,
      33 words.
    \item BPE builds the vocabulary by repeatedly merging the most frequent adjacent pair inside
      pre-tokenized words; the number of merges trades vocabulary size against length.
    \item Encoding replays the merges in order, so unseen words are spelled from known pieces.
    \item Starting from the 256 byte values makes nothing unknown. GPT-2: $256 + 50{,}000 + 1 =
      50{,}257$ tokens.
    \item ``The cat sat on the'' is five GPT-2 tokens, 464 3797 3332 319 262; the tokenizer is
      lossless, and capitals and spaces make different tokens.
    \item Tokenization shows through in spelling, digits and other languages, and token counts are
      what you pay for and what you run out of.
  \end{enumerate}
\end{frame}

\begin{frame}{Exercises: try it in the notebook}
  \small
  \begin{enumerate}
    \item \textbf{By hand.} On the toy text, find the most frequent adjacent pair inside the
      pre-tokenized words (each with its leading space) before running anything, and check your answer against merge 1 (cell In
      [\primercell{89}]). Why is the tie broken the way it is?
    \item \textbf{The dial} (cell In [\primercell{89}]). Train with \free{5 and then 30} merges.
      Predict each vocabulary size from the rule (characters plus one per merge) and say which run
      leaves the text shorter. When does the loop stop early?
    \item \textbf{Unseen words} (cell In [\primercell{165}]). First re-run cell In [\primercell{89}]
      with 10 merges, so \texttt{merges} is the guide's list again. Then encode \free{``the dogs ran''} and
      \free{``The cat''} with the toy encoder. Predict which pieces come out and which have no ID.
    \item \textbf{Bytes} (cell In [\primercell{196}]). Count the bytes of a word in a language you
      know that isn't English. How many tokens does GPT-2 give it (cell In [\primercell{261}])?
    \item \textbf{Your name} (cell In [\primercell{261}]). Encode your name with and without a
      leading space, and in capitals. Explain why the IDs differ.
  \end{enumerate}
\end{frame}

\begin{frame}{Further reading}
  \small
  \begin{itemize}
    \item P.\ Gage (1994). A new algorithm for data compression. \emph{C Users Journal}.
    \item R.\ Sennrich, B.\ Haddow \& A.\ Birch (2016). Neural machine translation of rare words with
      subword units. \emph{ACL}.
    \item A.\ Radford et al.\ (2019). Language models are unsupervised multitask learners.
      \emph{OpenAI}. \ (GPT-2)
    \item T.\ Kudo (2018). Subword regularization: improving neural network translation models with
      multiple subword candidates. \emph{ACL}. \ (Unigram)
    \item T.\ Kudo \& J.\ Richardson (2018). SentencePiece: a simple and language independent subword
      tokenizer and detokenizer for neural text processing. \emph{EMNLP}.
    \item A.\ Kaushal \& K.\ Mahowald (2022). What do tokens know about their characters and how do
      they know it? \emph{NAACL}.
  \end{itemize}
  {\small\color{primerMuted}Links for each, and for WordPiece, BERT and T5, are in the
  \href{\guideurl}{guide}.}
\end{frame}

\begin{frame}{Next time: Guide 3, embeddings}
  \begin{itemize}
    \item Today's output is a sequence of integers, and an integer is only a label.
    \item Guide 3 gives each token ID a learned vector the network can compute with, by letting
      each ID select one row of a table.
  \end{itemize}
  \vspace{4mm}
  \begin{keyidea}[Before next time]
    Run the notebook from top to bottom and check that every printed number matches the slides.
  \end{keyidea}
\end{frame}

\end{document}
